\documentclass[12pt,a4paper]{article}

\usepackage[utf8]{inputenc}
\usepackage[T1]{fontenc}
\usepackage{amsmath, amssymb}
\usepackage{enumitem}
\usepackage{geometry}
\usepackage{xcolor}
\usepackage{tcolorbox}
\usepackage{hyperref}

\geometry{margin=2.5cm}

\title{Banco de Questões -- Aula 09\\Dependência Funcional e Normalização}
\author{Banco de Dados}
\date{}

\newtcolorbox{resposta}{
  colback=gray!5,
  colframe=black!40,
  title=Resposta comentada
}

\begin{document}

\maketitle

\section*{Orientações}

Este banco de questões cobre os principais tópicos da Aula 09:
\begin{itemize}
    \item medidas informais de qualidade de esquemas relacionais;
    \item semântica dos atributos;
    \item redundância;
    \item valores nulos;
    \item tuplas ilegítimas;
    \item dependência funcional;
    \item formas normais;
    \item 1FN, 2FN, 3FN e BCNF;
    \item cobertura mínima de dependências funcionais.
\end{itemize}

\section{Questões Conceituais}

\begin{enumerate}[label=\textbf{Q\arabic*.}]

\item O que a normalização busca melhorar em um projeto de banco de dados relacional?

\begin{resposta}
A normalização busca organizar os atributos em relações de forma mais adequada, reduzindo redundâncias, evitando anomalias e melhorando a qualidade do esquema relacional. Ela ajuda a garantir que cada tabela represente bem um conjunto de fatos do minimundo.
\end{resposta}

\item Cite as quatro medidas informais de qualidade apresentadas na aula.

\begin{resposta}
As quatro medidas são:
\begin{enumerate}
    \item semântica dos atributos;
    \item redução de valores redundantes nas tuplas;
    \item redução de valores nulos nas tuplas;
    \item impedimento da geração de tuplas ilegítimas.
\end{enumerate}
\end{resposta}

\item O que significa dizer que uma relação deve ter boa semântica dos atributos?

\begin{resposta}
Significa que os atributos agrupados em uma mesma relação devem fazer sentido juntos no mundo real. Uma tabela deve representar um único conceito principal ou um conjunto coerente de fatos.
\end{resposta}

\item Por que a redundância é um problema em esquemas relacionais?

\begin{resposta}
Porque dados repetidos aumentam o espaço ocupado e podem causar inconsistências. Se uma informação aparece em várias tuplas, uma atualização incompleta pode deixar o banco com versões diferentes da mesma informação.
\end{resposta}

\item Explique o que é uma anomalia de atualização.

\begin{resposta}
É um problema que ocorre quando uma informação repetida precisa ser alterada em vários lugares. Se a alteração for feita apenas em algumas tuplas, o banco fica inconsistente.
\end{resposta}

\item Explique o que é uma anomalia de inserção.

\begin{resposta}
É um problema que ocorre quando não conseguimos inserir uma informação sem inserir outra informação desnecessária. Por exemplo, não conseguir cadastrar uma disciplina enquanto ainda não houver aluno matriculado nela.
\end{resposta}

\item Explique o que é uma anomalia de remoção.

\begin{resposta}
É um problema que ocorre quando, ao remover uma tupla, perdemos também uma informação que deveria continuar armazenada. Por exemplo, excluir o último aluno de uma disciplina e, com isso, perder os dados da disciplina.
\end{resposta}

\item Por que muitos valores \texttt{NULL} podem indicar um problema de projeto?

\begin{resposta}
Porque muitos valores nulos podem indicar que atributos de conceitos diferentes estão misturados na mesma relação. Isso pode dificultar consultas, restrições e interpretação dos dados.
\end{resposta}

\item O que é uma dependência funcional?

\begin{resposta}
Uma dependência funcional é uma restrição entre dois conjuntos de atributos. Dizemos que $X \rightarrow Y$ quando os valores de $X$ determinam os valores de $Y$. Ou seja, se duas tuplas têm o mesmo valor em $X$, então devem ter o mesmo valor em $Y$.
\end{resposta}

\item Interprete a dependência funcional:
\[
\text{matricula} \rightarrow \text{nomeAluno}
\]

\begin{resposta}
Significa que a matrícula determina o nome do aluno. Se duas tuplas possuem a mesma matrícula, elas devem possuir o mesmo nome de aluno.
\end{resposta}

\end{enumerate}

\section{Questões de Aplicação}

\begin{enumerate}[resume,label=\textbf{Q\arabic*.}]

\item Considere a relação:
\[
Aluno(\text{matricula}, \text{nome}, \text{curso}, \text{coordenador})
\]

Suponha que cada matrícula identifica um aluno e cada curso possui um único coordenador. Indique dependências funcionais possíveis.

\begin{resposta}
Algumas dependências funcionais possíveis são:
\[
\text{matricula} \rightarrow \text{nome}
\]
\[
\text{matricula} \rightarrow \text{curso}
\]
\[
\text{curso} \rightarrow \text{coordenador}
\]

Como a matrícula identifica o aluno, ela determina o nome e o curso. Como cada curso possui um único coordenador, o curso determina o coordenador.
\end{resposta}

\item Na relação da questão anterior, existe redundância? Justifique.

\begin{resposta}
Sim. O coordenador do curso pode aparecer repetido para vários alunos do mesmo curso. Se o coordenador mudar, será necessário atualizar várias tuplas, o que pode gerar anomalia de atualização.
\end{resposta}

\item A relação
\[
Aluno(\text{matricula}, \text{nome}, \text{curso}, \text{coordenador})
\]
poderia ser decomposta em quais relações para reduzir redundância?

\begin{resposta}
Uma decomposição possível é:
\[
Aluno(\text{matricula}, \text{nome}, \text{curso})
\]
\[
Curso(\text{curso}, \text{coordenador})
\]

Assim, os dados do curso ficam separados dos dados do aluno.
\end{resposta}

\item Considere:
\[
Pedido(\text{idPedido}, \text{idCliente}, \text{nomeCliente}, \text{dataPedido})
\]

Sabendo que cada cliente possui um único nome, indique uma possível dependência funcional problemática.

\begin{resposta}
A dependência problemática é:
\[
\text{idCliente} \rightarrow \text{nomeCliente}
\]

Se o nome do cliente fica armazenado em todos os pedidos, haverá redundância.
\end{resposta}

\item Por que a relação da questão anterior pode gerar anomalia de atualização?

\begin{resposta}
Porque o nome do cliente pode aparecer repetido em vários pedidos. Se o nome for corrigido ou alterado, será necessário atualizar todos os pedidos daquele cliente.
\end{resposta}

\item Considere a relação:
\[
Matricula(\text{aluno}, \text{disciplina}, \text{professor}, \text{sala})
\]

Suponha que:
\[
(\text{aluno}, \text{disciplina}) \rightarrow \text{professor}
\]
\[
\text{disciplina} \rightarrow \text{sala}
\]

Qual atributo depende apenas de parte de uma chave composta?

\begin{resposta}
Se a chave da relação for $(aluno, disciplina)$, então \texttt{sala} depende apenas de \texttt{disciplina}. Logo, existe uma dependência parcial:
\[
\text{disciplina} \rightarrow \text{sala}
\]
Isso viola a ideia da Segunda Forma Normal.
\end{resposta}

\item Explique por que a dependência
\[
\text{disciplina} \rightarrow \text{sala}
\]
pode causar redundância.

\begin{resposta}
Porque a sala da disciplina será repetida para todos os alunos matriculados naquela disciplina. Se a sala mudar, várias tuplas precisarão ser atualizadas.
\end{resposta}

\item Considere a relação:
\[
Funcionario(\text{cpf}, \text{nome}, \text{idDepto}, \text{nomeDepto})
\]

Com as dependências:
\[
\text{cpf} \rightarrow \text{nome}, \text{idDepto}
\]
\[
\text{idDepto} \rightarrow \text{nomeDepto}
\]

Existe dependência transitiva? Explique.

\begin{resposta}
Sim. Como:
\[
\text{cpf} \rightarrow \text{idDepto}
\]
e
\[
\text{idDepto} \rightarrow \text{nomeDepto}
\]
então:
\[
\text{cpf} \rightarrow \text{nomeDepto}
\]

Nesse caso, \texttt{nomeDepto} depende de \texttt{cpf} de forma transitiva, passando por \texttt{idDepto}.
\end{resposta}

\item A relação da questão anterior está adequada à Terceira Forma Normal? Justifique.

\begin{resposta}
Não. Há dependência transitiva de um atributo não-chave em relação à chave. O nome do departamento deveria ficar em uma relação separada, como:
\[
Funcionario(\text{cpf}, \text{nome}, \text{idDepto})
\]
\[
Departamento(\text{idDepto}, \text{nomeDepto})
\]
\end{resposta}

\item Explique, com suas palavras, a diferença entre dependência parcial e dependência transitiva.

\begin{resposta}
A dependência parcial ocorre quando um atributo depende apenas de parte de uma chave composta. A dependência transitiva ocorre quando um atributo depende da chave por meio de outro atributo não-chave.
\end{resposta}

\end{enumerate}

\section{Questões sobre Formas Normais}

\begin{enumerate}[resume,label=\textbf{Q\arabic*.}]

\item O que é a Primeira Forma Normal?

\begin{resposta}
Uma relação está na Primeira Forma Normal quando seus atributos possuem valores atômicos, ou seja, não há grupos repetidos ou atributos multivalorados dentro de uma mesma célula.
\end{resposta}

\item A relação abaixo viola a Primeira Forma Normal?
\[
Aluno(\text{matricula}, \text{nome}, \text{telefones})
\]

Exemplo de tupla:
\[
(123, \text{Iury}, \{9999-1111, 9888-2222\})
\]

\begin{resposta}
Sim. O atributo \texttt{telefones} possui múltiplos valores dentro da mesma célula. Para ficar em 1FN, os telefones deveriam ser separados, por exemplo:
\[
Aluno(\text{matricula}, \text{nome})
\]
\[
TelefoneAluno(\text{matricula}, \text{telefone})
\]
\end{resposta}

\item O que é a Segunda Forma Normal?

\begin{resposta}
Uma relação está na Segunda Forma Normal quando está na Primeira Forma Normal e não possui dependências parciais de atributos não-chave em relação a uma chave composta.
\end{resposta}

\item Uma relação com chave primária simples pode violar a Segunda Forma Normal por dependência parcial?

\begin{resposta}
Não. Dependência parcial só faz sentido quando a chave é composta. Se a chave possui apenas um atributo, não há ``parte'' da chave.
\end{resposta}

\item O que é a Terceira Forma Normal?

\begin{resposta}
Uma relação está na Terceira Forma Normal quando está na Segunda Forma Normal e não possui dependências transitivas de atributos não-chave em relação à chave.
\end{resposta}

\item O que é a Forma Normal de Boyce-Codd?

\begin{resposta}
A Forma Normal de Boyce-Codd, ou BCNF, é uma forma normal mais rigorosa que a 3FN. Em geral, uma relação está em BCNF quando, para toda dependência funcional não trivial $X \rightarrow Y$, o conjunto $X$ é uma superchave.
\end{resposta}

\item Toda relação em BCNF está em 3FN? E toda relação em 3FN está em BCNF?

\begin{resposta}
Toda relação em BCNF está em 3FN, pois BCNF é mais restritiva. Porém, nem toda relação em 3FN está em BCNF.
\end{resposta}

\item Considere:
\[
R(A, B, C)
\]
com dependências:
\[
A \rightarrow B
\]
\[
B \rightarrow C
\]

Se $A$ é chave, há violação de 3FN? Justifique.

\begin{resposta}
Sim, pode haver violação de 3FN por dependência transitiva. Como:
\[
A \rightarrow B
\]
e
\[
B \rightarrow C
\]
então:
\[
A \rightarrow C
\]
por meio de $B$. Se $B$ não é chave e $C$ é não-chave, temos dependência transitiva.
\end{resposta}

\item Considere:
\[
R(A, B, C)
\]
com dependências:
\[
A \rightarrow B
\]
\[
B \rightarrow A
\]
\[
B \rightarrow C
\]

Quais poderiam ser chaves candidatas?

\begin{resposta}
Como $A \rightarrow B$ e $B \rightarrow A$, $A$ e $B$ determinam um ao outro. Além disso, $B \rightarrow C$. Portanto, $B$ determina $A$ e $C$, então $B$ é chave candidata.

Como $A \rightarrow B$ e $B \rightarrow C$, então $A \rightarrow C$. Logo, $A$ também determina $B$ e $C$. Portanto, $A$ também é chave candidata.
\end{resposta}

\item Por que a normalização não deve ser vista apenas como uma regra mecânica?

\begin{resposta}
Porque a normalização depende do significado dos dados no minimundo. As dependências funcionais não surgem apenas olhando os nomes dos atributos; elas dependem das regras reais do domínio modelado.
\end{resposta}

\end{enumerate}

\section{Questões sobre Cobertura Mínima}

\begin{enumerate}[resume,label=\textbf{Q\arabic*.}]

\item O que é uma cobertura mínima de dependências funcionais?

\begin{resposta}
É um conjunto equivalente de dependências funcionais, mas sem redundâncias. Em uma cobertura mínima, normalmente cada dependência tem apenas um atributo no lado direito, não há atributos redundantes no lado esquerdo e não há dependências funcionais desnecessárias.
\end{resposta}

\item Qual é o primeiro passo comum para calcular uma cobertura mínima?

\begin{resposta}
Separar dependências com vários atributos no lado direito. Por exemplo:
\[
A \rightarrow BC
\]
deve ser substituída por:
\[
A \rightarrow B
\]
\[
A \rightarrow C
\]
\end{resposta}

\item Considere:
\[
F = \{A \rightarrow BC, B \rightarrow D\}
\]

Reescreva $F$ colocando apenas um atributo no lado direito de cada dependência.

\begin{resposta}
O conjunto fica:
\[
F = \{A \rightarrow B,\ A \rightarrow C,\ B \rightarrow D\}
\]
\end{resposta}

\item Considere:
\[
F = \{AB \rightarrow C,\ A \rightarrow C\}
\]

O atributo $B$ é redundante em $AB \rightarrow C$? Justifique.

\begin{resposta}
Sim. Como já existe $A \rightarrow C$, não precisamos de $B$ para determinar $C$. Logo, $AB \rightarrow C$ pode ser simplificada para $A \rightarrow C$.
\end{resposta}

\item Considere:
\[
F = \{A \rightarrow B,\ B \rightarrow C,\ A \rightarrow C\}
\]

A dependência $A \rightarrow C$ é redundante? Justifique.

\begin{resposta}
Sim. Como $A \rightarrow B$ e $B \rightarrow C$, podemos inferir $A \rightarrow C$ por transitividade. Logo, $A \rightarrow C$ pode ser removida sem alterar o fechamento do conjunto.
\end{resposta}

\end{enumerate}

\section{Questões Estilo Prova}

\begin{enumerate}[resume,label=\textbf{Q\arabic*.}]

\item Considere a relação:
\[
Venda(\text{idVenda}, \text{idProduto}, \text{nomeProduto}, \text{idCliente}, \text{nomeCliente}, \text{dataVenda})
\]

Suponha:
\[
\text{idVenda} \rightarrow \text{idProduto}, \text{idCliente}, \text{dataVenda}
\]
\[
\text{idProduto} \rightarrow \text{nomeProduto}
\]
\[
\text{idCliente} \rightarrow \text{nomeCliente}
\]

Identifique problemas de projeto nessa relação e proponha uma decomposição.

\begin{resposta}
A relação mistura informações da venda, do produto e do cliente. Há redundância em \texttt{nomeProduto} e \texttt{nomeCliente}, pois esses dados podem se repetir em várias vendas.

Uma decomposição possível é:
\[
Venda(\text{idVenda}, \text{idProduto}, \text{idCliente}, \text{dataVenda})
\]
\[
Produto(\text{idProduto}, \text{nomeProduto})
\]
\[
Cliente(\text{idCliente}, \text{nomeCliente})
\]

Assim, cada relação representa um conceito mais claro.
\end{resposta}

\item Considere:
\[
R(A, B, C, D)
\]
com as dependências:
\[
A \rightarrow B
\]
\[
B \rightarrow C
\]
\[
C \rightarrow D
\]

Determine o fechamento de $A$, isto é, $A^+$.

\begin{resposta}
Começamos com:
\[
A^+ = \{A\}
\]

Como $A \rightarrow B$:
\[
A^+ = \{A, B\}
\]

Como agora temos $B$ e $B \rightarrow C$:
\[
A^+ = \{A, B, C\}
\]

Como agora temos $C$ e $C \rightarrow D$:
\[
A^+ = \{A, B, C, D\}
\]

Logo:
\[
A^+ = \{A, B, C, D\}
\]

Portanto, $A$ é chave para $R$.
\end{resposta}

\item Considere:
\[
R(A, B, C, D)
\]
com:
\[
AB \rightarrow C
\]
\[
C \rightarrow D
\]

Calcule $(AB)^+$.

\begin{resposta}
Começamos com:
\[
(AB)^+ = \{A, B\}
\]

Como $AB \rightarrow C$:
\[
(AB)^+ = \{A, B, C\}
\]

Como $C \rightarrow D$:
\[
(AB)^+ = \{A, B, C, D\}
\]

Logo:
\[
(AB)^+ = \{A, B, C, D\}
\]

Portanto, $AB$ é chave para $R$.
\end{resposta}

\item Considere:
\[
R(A, B, C)
\]
com:
\[
A \rightarrow B
\]
\[
B \rightarrow A
\]
\[
B \rightarrow C
\]

A relação está em BCNF? Justifique.

\begin{resposta}
Para verificar BCNF, analisamos se o lado esquerdo de cada dependência funcional é superchave.

Como $B \rightarrow A$ e $B \rightarrow C$, então:
\[
B^+ = \{A, B, C\}
\]
Logo, $B$ é superchave.

Como $A \rightarrow B$ e $B \rightarrow C$, então:
\[
A^+ = \{A, B, C\}
\]
Logo, $A$ também é superchave.

Como o lado esquerdo das dependências relevantes é superchave, a relação está em BCNF.
\end{resposta}

\item Considere:
\[
R(A, B, C)
\]
com:
\[
A \rightarrow B
\]
\[
B \rightarrow C
\]

Sabendo que $A$ é chave, proponha uma decomposição para reduzir dependência transitiva.

\begin{resposta}
Como $A \rightarrow B$ e $B \rightarrow C$, temos uma dependência transitiva:
\[
A \rightarrow C
\]
por meio de $B$.

Uma decomposição possível é:
\[
R_1(A, B)
\]
\[
R_2(B, C)
\]

Assim, a dependência $B \rightarrow C$ fica em uma relação própria.
\end{resposta}

\end{enumerate}

\section*{Resumo para Revisão Rápida}

\begin{itemize}
    \item Dependência funcional: $X \rightarrow Y$ significa que $X$ determina $Y$.
    \item Redundância pode gerar anomalias de inserção, atualização e remoção.
    \item 1FN elimina atributos multivalorados ou não atômicos.
    \item 2FN elimina dependências parciais em relação a chaves compostas.
    \item 3FN elimina dependências transitivas de atributos não-chave.
    \item BCNF exige que, em toda DF não trivial $X \rightarrow Y$, $X$ seja superchave.
    \item Cobertura mínima remove redundâncias do conjunto de dependências funcionais.
\end{itemize}

\end{document}
